\section{The simply typed lambda calculus}
\label{sec:normalization-curry-howard-02}

Fix a set \(\mathsf{Base}\) of base types.  Types and terms are generated by
the following scoped metavariables: in this section only, \(A,B\) range over
types, \(M,N\) range over terms, and \(\alpha\) ranges over members of
\(\mathsf{Base}\).  The symbol \(\Theta_{\mathrm{ty}}\) ranges over finite
lists assigning distinct term variables to types; these lists are the typing
contexts of this section.  These are metalinguistic roles, not additional
object-language variables.
\[
  A,B::=\alpha\mid A\to B,
  \qquad \alpha\in\mathsf{Base},
\]
and
\[
  M,N::=x\mid \lambda x{:}A.M\mid MN.
\]
The typing judgment \(\Theta_{\mathrm{ty}}\derives M:A\) is generated by
\[
\frac{x{:}A\in\Theta_{\mathrm{ty}}}
     {\Theta_{\mathrm{ty}}\derives x:A}\;\textsc{Var},
\qquad
\frac{\Theta_{\mathrm{ty}},x{:}A\derives M:B}
     {\Theta_{\mathrm{ty}}\derives\lambda x{:}A.M:A\to B}
     \;\textsc{\(\to I\)},
\]
In \textsc{\(\to I\)}, the variable \(x\) does not occur in
\(\Theta_{\mathrm{ty}}\); this can always be arranged by renaming the
bound variable of the abstraction, and typing judgments are stable under
such renaming of the subject term.
\[
\frac{\Theta_{\mathrm{ty}}\derives M:A\to B
      \qquad
      \Theta_{\mathrm{ty}}\derives N:A}
     {\Theta_{\mathrm{ty}}\derives MN:B}\;\textsc{\(\to E\)}.
\]
The principal computation is beta reduction:
\[
  (\lambda x{:}A.M)N\longrightarrow_\beta M[N/x],
\]
where term substitution is capture-avoiding.

\begin{theorem}[Subject reduction]
If \(\Theta_{\mathrm{ty}}\derives M:A\) and
\(M\longrightarrow_\beta N\), then
\(\Theta_{\mathrm{ty}}\derives N:A\).
\end{theorem}

\paragraph{Proof sketch.}
Prove the typing substitution lemma: if
\(\Theta_{\mathrm{ty}},x{:}B\derives M:A\) and
\(\Theta_{\mathrm{ty}}\derives N:B\), then
\(\Theta_{\mathrm{ty}}\derives M[N/x]:A\).  Apply it to the contracted beta
redex and use induction on the compatible closure of beta reduction for
reductions inside a term.

\begin{theorem}[Strong normalization]
Every well-typed term in the simply typed lambda calculus has no infinite
beta-reduction sequence.
\end{theorem}

\paragraph{Proof sketch.}
Assign to each type a reducibility candidate.  A term belongs to the candidate
for a base type when it is strongly normalizing; it belongs to the candidate
for \(A\to B\) when application to every member of the candidate for \(A\)
produces a member of the candidate for \(B\).  The fundamental lemma proves by
induction on typing derivations that every well-typed term belongs to its
type's candidate.  Strong normalization follows.  This proof sketch gives the
logical-relations architecture; closure lemmas for the candidates are omitted.

